% Part: methods
% Chapter: proofs
% Section: example-1

\documentclass[../../../include/open-logic-section]{subfiles}

\begin{document}

\olfileid{mth}{prf}{ex1}

\olsection{एक उदाहरण}

पहिले उदाहरण म्हणजे संचांच्या संयोग आणि छेदाविषयीचे पुढील साधे तथ्य. त्यातून व्याख्या उलगडणे, संधियोग आणि वैश्विक विधाने सिद्ध करणे, तसेच प्रकरणांनुसार सिद्धता करणे या पद्धती स्पष्ट होतील.

\begin{prop}
कोणत्याही $A$, $B$ आणि $C$ संचांसाठी, $A \cup (B \cap C) = (A \cup B) \cap (A \cup C)$.
\end{prop}

हे सिद्ध करू!{}

\begin{proof}
कोणत्याही $A$, $B$ आणि $C$ संचांसाठी $A \cup (B \cap C) = (A \cup B) \cap (A \cup C)$ दाखवायचे आहे.
\begin{quote}
प्रथम विधानातील “$=$” ची व्याख्या उलगडू. संच एकरूप आहेत हे सिद्ध करणे म्हणजे त्यांचे !!{element}s तेच आहेत हे दाखवणे, हे आठवा. म्हणजे $A \cup (B \cap C)$ मधील सर्व !!{element}s हे $(A \cup B) \cap (A \cup C)$ मधील !!{element}s देखील असले पाहिजेत आणि उलटही तसेच असले पाहिजे. “उलटही तसेच” म्हणजे $(A \cup B) \cap (A \cup C)$ मधील प्रत्येक !!{element} हा $A \cup (B \cap C)$ मधील !!a{element} देखील असला पाहिजे. म्हणून व्याख्या उलगडल्यावर संधियोग सिद्ध करावा लागतो असे दिसते. हे नोंदवू:
\end{quote}
व्याख्येनुसार, $A \cup (B \cap C) = (A \cup B) \cap (A \cup C)$ हे तेव्हाच आणि फक्त तेव्हा सत्य असते जेव्हा $A \cup (B \cap C)$ मधील प्रत्येक !!{element} हा $(A \cup B) \cap (A \cup C)$ मधील !!a{element} देखील असतो आणि $(A \cup B) \cap (A \cup C)$ मधील प्रत्येक !!{element} हा $A \cup (B \cap C)$ मधील !!a{element} असतो.
\begin{quote}
हा संधियोग असल्यामुळे प्रत्येक संधिघटक वेगळा सिद्ध करावा लागेल. पहिल्यापासून सुरुवात करू: $A \cup (B \cap C)$ मधील प्रत्येक !!{element} हा $(A \cup B) \cap (A \cup C)$ मधील !!a{element} देखील असतो हे सिद्ध करू.

हे वैश्विक विधान आहे. म्हणून $A \cup (B \cap C)$ मधील एक स्वेच्छ !!{element} घेऊन तो $(A \cup B) \cap (A \cup C)$ मधील !!a{element} देखील असला पाहिजे हे दाखवू. या स्वेच्छ !!{element} साठी एखादे चल निवडू; उदाहरणार्थ,~$z$. सिद्धता पुढे अशी चालते:
\end{quote}
प्रथम $A \cup (B \cap C)$ मधील प्रत्येक !!{element} हा $(A \cup B) \cap (A \cup C)$ मधील !!a{element} देखील असतो हे सिद्ध करू. $z \in A \cup (B \cap C)$ असू द्या. $z \in (A \cup B) \cap (A \cup C)$ दाखवायचे आहे.
\begin{quote}
आता $\cup$ आणि~$\cap$ यांच्या व्याख्या उलगडायच्या आहेत. उदाहरणार्थ, $\cup$ ची व्याख्या अशी आहे: $A \cup B = \Setabs{z}{z \in A \text{ किंवा } z \in B}$. “$A \cup (B \cap C)$” ला ही व्याख्या लावताना व्याख्येतील “$B$” ची भूमिका आता “$B \cap C$” बजावते. म्हणून $A \cup (B \cap C) = \Setabs{z}{z \in A \text{ किंवा } z \in B \cap C}$. त्यामुळे $z \in A \cup (B \cap C)$ या गृहीतकाचा अर्थ $z \in \Setabs{z}{z \in A \text{ किंवा } z \in B \cap C}$ असा होतो. तसेच $z \in \Setabs{z}{\dots z\dots}$ हे \dots $z$ \dots{} असते तेव्हाच आणि फक्त तेव्हाच सत्य असते. म्हणजे येथे $z \in A$ किंवा $z \in B \cap C$ असते.
\end{quote}
$\cup$ च्या व्याख्येनुसार $z \in A$ किंवा $z \in B \cap C$ असते.
\begin{quote}
हा विकल्पयोग असल्यामुळे प्रकरणांनुसार सिद्धता उपयुक्त ठरेल. दोन्ही प्रकरणे घेऊन प्रत्येक प्रकरणात अपेक्षित निष्कर्ष---म्हणजे “$z \in (A \cup B) \cap (A \cup C)$”---मिळतो हे दाखवू.
\end{quote}
प्रकरण 1: $z \in A$ असे समजा.
\begin{quote}
या गृहीतकांवरून आणखी फार काही मिळत नाही. म्हणून अपेक्षित निष्कर्षात कोणती रचना आहे ते पाहू. $z \in (A \cup B) \cap (A \cup C)$ दाखवायचे आहे. $\cap$ च्या व्याख्येनुसार $z \in (A \cup B) \cap (A \cup C)$ दाखवण्यासाठी तो $(A \cup B)$ आणि $(A \cup C)$ या दोन्हींत आहे हे दाखवावे लागते. पण $z \in A \cup B$ हे $z \in A$ किंवा $z \in B$ असते तेव्हाच आणि फक्त तेव्हाच सत्य असते. तसेच प्रकरण~1 चे गृहीतक $z \in A$ आधीच उपलब्ध आहे. हाच युक्तिवाद $B$ ऐवजी $C$ वापरून केल्यास $z \in A \cup C$ मिळते. हा विचार अपेक्षित निष्कर्षापासून मागे जात केलेला होता; आता सिद्धतेत आवश्यक असलेल्या पुढच्या दिशेने तो नोंदवू.
\end{quote}
$z \in A$ असल्यामुळे $z \in A$ किंवा $z \in B$. म्हणून~$\cup$ च्या व्याख्येनुसार $z \in A \cup B$. त्याचप्रमाणे $z \in A \cup C$. पण~$\cap$ च्या व्याख्येनुसार याचाच अर्थ $z \in (A \cup B) \cap (A \cup C)$ असा होतो.
\begin{quote}
यामुळे प्रकरणांनुसार सिद्धतेचे पहिले प्रकरण पूर्ण झाले. आता $z \in B \cap C$ असलेल्या दुसऱ्या प्रकरणात निष्कर्ष मिळवायचा आहे.
\end{quote}
प्रकरण 2: $z \in B \cap C$ असे समजा.
\begin{quote}
पुन्हा दोन संचांचा छेद हाताळत आहोत. ~$ \cap$ ची व्याख्या वापरू:
\end{quote}
$z \in B \cap C$ असल्यामुळे~$\cap$ च्या व्याख्येनुसार $z$ हा $B$ आणि $C$ या दोन्हींचा !!a{element} असला पाहिजे.
\begin{quote}
पुन्हा अपेक्षित निष्कर्ष पाहू. $z$ हा $(A \cup B)$ आणि $(A \cup C)$ या दोन्हींत आहे हे दाखवायचे आहे. पुन्हा उत्तर लगेच मिळते.
\end{quote}
$z \in B$ असल्यामुळे $z \in (A \cup B)$. $z \in C$ असल्यामुळे $z \in (A \cup C)$ देखील. म्हणून $z \in (A \cup B) \cap (A \cup C)$.
\begin{quote}
येथे पुन्हा $\cup$ आणि $\cap$ यांच्या व्याख्या वापरल्या. पण त्या व्याख्या आधीच आठवून घेतल्या आहेत आणि $z$ हा दोन संचांपैकी एका संचात असेल तर त्यांच्या संयोगात असतो, हेही आधीच दाखवले आहे. म्हणून आता प्रत्येक वापराचा तितका तपशील लिहिण्याची गरज नाही.

प्रकरणांनुसार सिद्धतेचे दुसरे प्रकरण पूर्ण झाले. म्हणून पहिला अपेक्षित निष्कर्ष आता सांगता येतो.
\end{quote}
म्हणून $z \in A \cup (B \cap C)$ असेल, तर $z \in (A \cup B) \cap (A \cup C)$.
\begin{quote}
आता उलटी दिशा दाखवायची आहे: $(A \cup B) \cap (A \cup C)$ मधील प्रत्येक !!{element} हा $A \cup (B \cap C)$ मधील !!a{element} आहे. आधीप्रमाणे, पहिल्या संचातील एक स्वेच्छ !!{element} घेतला आहे असे गृहीत धरून तो दुसऱ्या संचात असला पाहिजे हे दाखवून हे वैश्विक विधान सिद्ध करू. काय करणार आहोत ते मांडू.
\end{quote}
आता $z \in (A \cup B) \cap (A \cup C)$ असे गृहीत धरा. $z \in A \cup (B \cap C)$ दाखवायचे आहे.
\begin{quote}
आता $z \in (A \cup B) \cap (A \cup C)$ या गृहीतकावरून विचार करत आहोत. सिद्धतेच्या पहिल्या भागात वापरलेले~$z$ हेच येथे वापरल्यामुळे फार गोंधळ होत नाही अशी आशा आहे. तो भाग पूर्ण झाल्यावर त्या भागातील तात्पुरती गृहीतके लागू राहिली नाहीत; म्हणून आता $z$ विषयी नवी गृहीतके घेता येतात. हे गोंधळात टाकत असेल, तर पुढील भागात $z$ ऐवजी वेगळे चल वापरा.

~$\cap$ च्या व्याख्येनुसार $z$ हा $A \cup B$ आणि $A \cup C$ या दोन्हींत आहे. पुढे $\cup$ ची व्याख्या उलगडल्यावर असे मिळते: $z \in A$ किंवा $z \in B$, आणि $z \in A$ किंवा $z \in C$ देखील. पुन्हा प्रकरणांनुसार सिद्धता दिसते; पण “आणि” मुळे गोंधळ होऊ शकतो. दोन्ही अटींऐवजी एकच अट पुरेशी आहे---$z$ हा $A$, $B$ किंवा $C$ पैकी एखाद्यात आहे---असे वाटू शकते. पण ती अदलाबदल चुकीची ठरेल. म्हणून काळजीपूर्वक एकेक विकल्पयोग पाहू.
\end{quote}
$\cap$ च्या व्याख्येनुसार $z \in A \cup B$ आणि $z \in A \cup C$. $\cup$ च्या व्याख्येनुसार $z \in A$ किंवा $z \in B$. प्रकरणे वेगळी पाहू.
\begin{quote}
पहिल्या विकल्पयोगावर लक्ष केंद्रित केल्यामुळे $A \cup C$ उलगडून मिळणारा दुसरा विकल्पयोग अजून नोंदवलेला नाही. आत्ता त्याची गरजही नाही. पहिले प्रकरण $z \in A$ हे आहे. संचाचा !!a{element} हा त्या संचाचा दुसऱ्या कोणत्याही संचाशी संयोग केल्यावर मिळणाऱ्या संचाचा !!a{element} देखील असतो. म्हणून प्रकरण~1 सोपे आहे:
\end{quote}
प्रकरण 1: $z \in A$ असे समजा. यावरून $z \in A \cup (B \cap C)$ मिळते.
\begin{quote}
आता दुसरे प्रकरण, $z \in B$, पाहू. येथे दुसऱ्या $\cup$ ची व्याख्या उलगडून आणखी एक प्रकरणांनुसार सिद्धता करू:
\end{quote}
प्रकरण 2: $z \in B$ असे समजा. $z \in A \cup C$ असल्यामुळे $z \in A$ किंवा $z \in C$. आणखी प्रकरणे वेगळी पाहू:

प्रकरण 2a: $z \in A$. मग पुन्हा $z \in A \cup (B \cap C)$.
\begin{quote}
हे थोडे विचित्र वाटेल: या प्रकरणात~$z \in B$ हे गृहीतक वापरण्याची गरजच पडली नाही. पण त्यात काही अडचण नाही.
\end{quote}
प्रकरण 2b: $z \in C$. मग $z \in B$ आणि $z \in C$. म्हणून $z \in B \cap C$ आणि परिणामी $z \in A \cup (B \cap C)$.
\begin{quote}
दोन्ही प्रकरणांनुसार सिद्धता पूर्ण झाल्या; त्यामुळे दुसरा अर्धा भागही पूर्ण झाला.
\end{quote}
म्हणून $z \in (A \cup B) \cap (A \cup C)$ असेल, तर $z \in A \cup (B \cap C)$.
\end{proof}

\end{document}
